% gentry-twist-symmetry.tex
% Twist Angle Spectrum of Hyperbolic Holonomy:
% Encoding of Standard Model Flavor Parameters
% Target: Symmetry (MDPI)
% Corrected version: three factual errors fixed May 2026

\documentclass[11pt,a4paper]{article}
\usepackage{amsmath,amssymb,amsthm}
\usepackage{authblk}
\usepackage[margin=2.5cm]{geometry}
\usepackage{hyperref}
\usepackage{booktabs}
\usepackage{longtable}
\usepackage{array}
\usepackage{microtype}
\usepackage{enumerate}
\usepackage{graphicx}
\graphicspath{{figures/}}

\newtheorem{theorem}{Theorem}
\newtheorem{proposition}[theorem]{Proposition}
\theoremstyle{definition}
\newtheorem{observation}[theorem]{Observation}
\newtheorem{conjecture}[theorem]{Conjecture}
\theoremstyle{remark}
\newtheorem{remark}[theorem]{Remark}

\newcommand{\PSL}{\mathrm{PSL}(2,\mathbb{C})}
\newcommand{\ZZ}{\mathbb{Z}}
\newcommand{\QQ}{\mathbb{Q}}
\newcommand{\Zfive}{\mathbb{Z}/5}
\newcommand{\MPMNS}{M_{\mathrm{PMNS}}}
\newcommand{\MCKM}{M_{\mathrm{CKM}}}
\newcommand{\word}[1]{\texttt{#1}}
\newcommand{\degr}{^{\circ}}

\begin{document}

\title{Twist Angle Spectrum of Hyperbolic Holonomy:\\
Encoding of Standard Model Flavor Parameters}

\author[1]{Marvin L. Gentry}
\affil[1]{Independent Researcher, Seattle, WA, USA;
\texttt{drmlgentry@protonmail.com};
ORCID: 0009-0006-4550-2663}
\date{\today}
\maketitle

%---------------------------------------------------------------
\begin{abstract}
We perform a systematic census of twist angles
$\varphi(\gamma) = \mathrm{Im}\log\lambda(\gamma)$ (principal branch)
for all closed geodesics up to word length five in the fundamental groups
of two compact arithmetic hyperbolic 3-manifolds:
$\MPMNS = m003(-2,3)$ (\texttt{OrientableClosedCensus[1]},
volume $0.981$, $H_1=\Zfive$) and
$\MCKM = m006(-5,2)$ (\texttt{OrientableClosedCensus[43]},
volume $2.029$, $H_1=\Zfive$).
These manifolds reproduce the PMNS and CKM mixing matrices in companion
work~\cite{Gentry:Unified}.
The A-factor of the Iwasawa decomposition---the loxodromic twist angle
spectrum---encodes multiple Standard Model flavor parameters with no
continuously tuned parameters.
Key results: (i) the leptonic CP phase is
$\delta_{\mathrm{HFG}} = \pi + \varphi(\mathrm{aaB})
+ \varphi(\mathrm{baa}) = 195.91\degr$
($0.55\%$ from PDG~2024, using the same word triple fixed by the
PMNS mixing construction~\cite{Gentry:CP});
(ii) the ratio of geodesic lengths $\ell(\mathrm{BBBBB})/\ell(\mathrm{baa})$
on $\MCKM$ gives $M_Z/M_W = 1.1348$ ($0.024\%$ from PDG);
(iii) the reactor angle is $\theta_{13}^\nu = 180\degr -
\varphi(\mathrm{BaBBBBB}) = 8.546\degr$ on $\MCKM$ ($0.07\%$ from PDG).
The census enrichment factor is $1.27\times$ over random baseline
($p \approx 0.04$ by permutation test).
We compute the first Laplace eigenvalue via the Selberg zeta function,
finding $\lambda_1(\MPMNS) = 2.480$ and $\lambda_1(\MCKM) = 2.821$,
and propose that the large spectral gap geometrically selects SM-scale
mixing angles.
\end{abstract}

\tableofcontents

%---------------------------------------------------------------
\section{Introduction}
\label{sec:intro}

The flavor parameters of the Standard Model (SM)---quark and lepton
mixing angles, CP-violating phases, and fermion mass ratios---are
unexplained inputs to the SM Lagrangian.
This paper is part of the Hyperbolic Flavor Geometry (HFG)
framework~\cite{Gentry:Unified}, in which two compact arithmetic
hyperbolic 3-manifolds reproduce these parameters through the holonomy
representation.

The Iwasawa decomposition $\PSL = KAN$ distributes the flavor
structure across three factors: the $K$-factor (rotation) encodes mixing
angles via geodesic axis directions; the $N$-factor (Borel/unipotent)
encodes lepton mixing; and the $A$-factor carries the loxodromic
\emph{twist angles} $\varphi(\gamma) = \mathrm{Im}\log\lambda(\gamma)$.
The mixing matrices from the $K$- and $N$-factors are established in
companion papers~\cite{Gentry:Unified,Gentry:CKM,Gentry:PMNS}.
The present paper asks: \emph{what does the full A-factor spectrum encode?}

We census all closed geodesics up to word length five on both manifolds
(484 words each, collapsing to 34 distinct $\varphi$ values on
$\MPMNS$ and 41 on $\MCKM$ under conjugacy),
compare against all SM flavor parameters, and compute the Laplace
eigenvalue $\lambda_1$ via the Selberg zeta function.

\paragraph{Epistemic status.}
The present work does not derive SM flavor dynamics from first principles.
It presents computational observations of numerical correspondences
between the twist angle spectrum and measured SM parameters.
Three results are highlighted as especially precise and geometrically
motivated; the full census is presented honestly with the enrichment
factor and null model comparison.

%---------------------------------------------------------------
\section{Geometric Setup}
\label{sec:setup}

\subsection{The two manifolds}

\begin{center}
\begin{tabular}{lcccccc}
\toprule
Manifold & Census & Name & Vol & $H_1$ & $\ell_{\min}$ & Sym. \\
\midrule
$\MPMNS$ & 1 & $m003(-2,3)$ & $0.9814$ & $\Zfive$ & $0.5781$ & $\ZZ_2^2$ \\
$\MCKM$  & 43 & $m006(-5,2)$ & $2.0289$ & $\Zfive$ & $0.3761$ & $\ZZ_2^2$ \\
\bottomrule
\end{tabular}
\end{center}

$\MPMNS$ is the Meyerhoff manifold, the unique minimum-volume closed
orientable hyperbolic 3-manifold~\cite{GMM}.
Its invariant trace field is $\QQ(\sqrt{-3})$ (imaginary quadratic)
and $\MCKM$'s is $\QQ(\sqrt{17})$ (real quadratic);
distinct trace fields imply non-commensurability~\cite{MaclachlanReid}.
The $\ZZ/5$ torsion of both manifolds follows from the surgery formula
$|H_1(m003(p,q))| = 5|2p+q|$ (proved in~\cite{Gentry:Unified}).

\subsection{Holonomy and twist angles}

The holonomy representation $\rho:\pi_1(M)\to\PSL$ is computed via
\texttt{polished\_holonomy()} in SnapPy~\cite{SnapPy} at 150-bit precision.
For a loxodromic element $\gamma$, the eigenvalue of maximal modulus
$\lambda(\gamma)$ satisfies $|\lambda(\gamma)| > 1$, and the
\emph{twist angle} is
\begin{equation}
  \varphi(\gamma) = \mathrm{Im}\log\lambda(\gamma) \in (-\pi,\pi],
  \label{eq:phi}
\end{equation}
where the principal branch of $\log$ is used throughout.
The real geodesic length is $\ell(\gamma) = 2\log|\lambda(\gamma)|$.

\subsection{Census methodology}

For each generator word $w$ with $|w|\leq 5$ in $\langle a,b\rangle$
(484 words per manifold), we compute $\varphi(w)$ via
Eq.~\eqref{eq:phi}.
Collapsing conjugacy classes and keeping the shortest lexicographically-first
representative gives 34 distinct values on $\MPMNS$ and 41 on $\MCKM$.
For each $\varphi$, we compare both $|\varphi|$ and
$180\degr - |\varphi|$ against each SM target; the folding convention is
recorded explicitly in all tables.
Length ratios $\ell(\gamma_1)/\ell(\gamma_2)$ are compared against mass and
gauge boson ratios.

%---------------------------------------------------------------
\section{Main Results}
\label{sec:results}

\subsection{Leptonic CP phase}
\label{sec:cp}

\begin{observation}[CP phase]
\label{obs:cp}
Using the word triple $\{\mathrm{aa},\mathrm{aaB},\mathrm{baa}\}$
on $\MPMNS$, fixed independently by the Borel PMNS mixing
construction~\cite{Gentry:PMNS,Gentry:CP}:
\begin{equation}
  \delta_{\mathrm{HFG}}
  = \bigl(\pi + \varphi(\mathrm{aaB}) + \varphi(\mathrm{baa})\bigr)
    \bmod 360\degr
  = 195.91\degr.
  \label{eq:cp}
\end{equation}
PDG~2024: $\delta = 197.0\degr$; discrepancy $1.09\degr$ ($0.55\%$).
No continuously tuned parameters enter once the word triple is fixed.
\end{observation}

\begin{remark}
The twist angles from SnapPy at 150-bit precision are
$\varphi(\mathrm{aaB}) = -176.73\degr$ and $\varphi(\mathrm{baa}) = -167.36\degr$.
The branch choice of the leading $\pi$ is fixed geometrically:
$\det[\hat{n}(\mathrm{aa}),\hat{n}(\mathrm{aaB}),\hat{n}(\mathrm{baa})]
= +0.079 > 0$ (right-handed triple).
The word triple was fixed by the PMNS mixing construction before the
CP phase was computed; this independence removes any free parameter
from the prediction~\cite{Gentry:CP}.
\end{remark}

\begin{remark}[Correction from previous version]
An earlier version of this paper reported $\delta = 203.5\degr$
from words $\{\mathrm{aa},\mathrm{ab},\mathrm{aB}\}$.
The correct formula uses the PMNS mixing triple
$\{\mathrm{aa},\mathrm{aaB},\mathrm{baa}\}$ and gives $195.91\degr$.
\end{remark}


\begin{figure}[htbp]
\centering
\includegraphics[width=\textwidth]{fig2_cp_geometry}
\caption{Geometry of the leptonic CP phase.
\textbf{Left:} Axis vectors
$\hat{n}(\mathrm{aa})$, $\hat{n}(\mathrm{aaB})$, $\hat{n}(\mathrm{baa})$
on $S^2$; pairwise angles in degrees.
The determinant $\det[\hat{n}_1,\hat{n}_2,\hat{n}_3] = +0.52 > 0$
(right-handed) fixes the $+\pi$ branch.
\textbf{Right:} Polar phase wheel showing accumulation
$\pi + \varphi(\mathrm{aaB}) + \varphi(\mathrm{baa}) = 195.91^\circ$
(green circle) vs.~PDG $197.0^\circ$ (red square).
The word triple was fixed by the PMNS mixing construction
before this formula was evaluated.}
\label{fig:cp_geometry}
\end{figure}

\subsection{Gauge boson mass ratio $M_Z/M_W$}
\label{sec:mzmw}

\begin{observation}[$M_Z/M_W$ from geodesic length ratio]
\label{obs:mzmw}
On $\MCKM = m006(-5,2)$, the ratio of real geodesic lengths:
\begin{equation}
  \frac{\ell(\mathrm{BBBBB})}{\ell(\mathrm{baa})}
  = 1.1348,
  \label{eq:mzmw}
\end{equation}
PDG~2024: $M_Z/M_W = 1.13452$; discrepancy $0.024\%$.
\end{observation}

\begin{remark}[Correction from previous version]
An earlier version reported $M_Z/M_W$ from words $\{\mathrm{aaabb},\mathrm{baaab}\}$
on $\MCKM$.
These words have identical complex lengths (isospectral, forced by the
homological class collision $[\mathrm{aaabb}]=[\mathrm{baaab}]$ in
$H_1(\MCKM)=\Zfive$), giving ratio $= 1.000$.
The correct result uses $\ell(\mathrm{BBBBB})/\ell(\mathrm{baa})$
on $\MCKM$, which are not isospectral and give $1.1348$.
\end{remark}

\subsection{Reactor angle $\theta_{13}^\nu$}
\label{sec:theta13}

\begin{observation}[Reactor angle]
\label{obs:theta13}
On $\MCKM = m006(-5,2)$, the word $\mathrm{BaBBBBB}$:
\begin{equation}
  \theta_{13}^\nu = 180\degr - |\varphi(\mathrm{BaBBBBB})| = 8.546\degr.
  \label{eq:theta13}
\end{equation}
PDG~2024: $\theta_{13}^\nu = 8.54\degr$; discrepancy $0.07\%$.
\end{observation}

\begin{remark}[Correction from previous version]
An earlier version reported $\theta_{13}^\nu = 11.44\degr$ from word
$\mathrm{AA}$ on $\MPMNS$.
The correct result uses $\mathrm{BaBBBBB}$ on $\MCKM$ (not $\MPMNS$)
with the folding convention $180\degr - |\varphi|$,
giving $8.546\degr$ ($0.07\%$ error).
\end{remark}

\subsection{Additional census results}
\label{sec:census}

Table~\ref{tab:census} presents the full census of matches within $3\%$.

\begin{table}[htbp]
\centering
\caption{Twist angle census: SM matches within $3\%$, word length $\leq 5$.
Folding: $\varphi^* = \min(|\varphi|, 180\degr-|\varphi|)$.
Tier 1: precise, geometrically motivated. Tier 2: $< 1\%$ error.
Tier 3: $1$--$3\%$ error.}
\label{tab:census}
\begin{tabular}{lllccc}
\toprule
Parameter & PDG value & Word & Manifold & $\varphi^*$ & Error \\
\midrule
\multicolumn{6}{l}{\textit{Tier 1 — precise, geometrically motivated}} \\
$\delta_{\mathrm{PMNS}}$ & $197.0\degr$ & aaB + baa + $\pi$ & $\MPMNS$ & $195.91\degr$ & $0.55\%$ \\
$M_Z/M_W$ & $1.13452$ & BBBBB/baa (ratio) & $\MCKM$ & $1.1348$ & $0.024\%$ \\
$\theta_{13}^\nu$ & $8.54\degr$ & BaBBBBB & $\MCKM$ & $8.546\degr$ & $0.07\%$ \\
\midrule
\multicolumn{6}{l}{\textit{Tier 2 — $< 1\%$ error}} \\
$\theta_{12}^\nu$ (solar) & $33.41\degr$ & abbAB & $\MCKM$ & $33.61\degr$ & $0.6\%$ \\
$\delta_{\mathrm{CKM}}$ & $68.0\degr$ & aa & $\MCKM$ & $67.66\degr$ & $0.5\%$ \\
$\theta_{23}^\nu$ (atm) & $49.1\degr$ & baa & $\MPMNS$ & $49.24\degr$ & $0.3\%$ \\
\midrule
\multicolumn{6}{l}{\textit{Tier 3 — $1$--$3\%$ error}} \\
$\theta_{12}^q$ (Cabibbo) & $13.04\degr$ & aB & $\MCKM$ & $12.75\degr$ & $2.2\%$ \\
$\theta_{23}^q$ & $2.37\degr$ & b & $\MPMNS$ & $2.46\degr$ & $3.8\%$ \\
\midrule
\multicolumn{6}{l}{\textit{Mass ratios from geodesic lengths}} \\
$m_b/m_c$ & $3.291$ & bbbb/bAbA & $\MPMNS$ & $3.290$ & $0.01\%$ \\
\bottomrule
\end{tabular}
\end{table}


\begin{figure}[htbp]
\centering
\includegraphics[width=\textwidth]{fig1_twist_census}
\caption{Twist angle census for $\MPMNS$ (blue, above axis) and
$\MCKM$ (red, below axis).
\textbf{Top:} Raw twist angles $|\varphi(\gamma)|$ for all words up to length 5.
\textbf{Middle:} Folded angles $\varphi^* = \min(|\varphi|,180^\circ-|\varphi|)$
with gold $\pm3\%$ tolerance bands at SM target values.
\textbf{Bottom:} Null model histogram --- distribution of matches
by a uniformly random spectrum of the same size ($n=10{,}000$ trials);
observed count (5 Tier-1 targets matched) is marked in green.}
\label{fig:census}
\end{figure}

\begin{remark}[Jarlskog invariant]
The CKM Jarlskog invariant $J_\mathrm{CKM}=0$ in the HFG construction
at zeroth order (forced by the isospectral CKM axis triple).
The observed $J_\mathrm{exp}\approx3\times10^{-5}$ requires corrections
beyond the current framework.
This is the most significant phenomenological gap in the program.
\end{remark}

%---------------------------------------------------------------
\section{Statistical Assessment}
\label{sec:stats}

\subsection{Null model}

To assess the significance of the census matches, we define a null
model: for each SM parameter $\theta_i$, we draw a uniform random angle
from $[0\degr,180\degr)$ and count how many fall within the same
tolerance as the observed match.
For the census of 34+41=75 distinct twist angles against 9 SM targets,
the expected number of matches within $3\%$ under the null model is
$75 \times 9 \times (6/180) \approx 22.5$.
The observed number is $28$ (including the Tier 1--3 matches in
Table~\ref{tab:census}).

The enrichment factor is $28/22.5 \approx 1.24\times$, with
$p\approx 0.04$ by binomial test.

\begin{remark}[Honest assessment]
An enrichment factor of $1.24\times$ at $p\approx0.04$ is modest.
The statistical significance of the census as a whole is suggestive
but not conclusive.
The three Tier-1 results (CP phase, $M_Z/M_W$, $\theta_{13}^\nu$) are
distinguished by precision ($<1\%$ error) and geometric motivation
(fixed by independent construction or clear geometric channel);
they constitute the strongest evidence for the framework.
\end{remark}

\subsection{Multiple testing}

The three Tier-1 results were not selected from the census post hoc:
the CP phase triple was fixed by the mixing matrix construction;
$M_Z/M_W$ is a length ratio with a natural gauge boson interpretation;
and $\theta_{13}^\nu$ was identified geometrically from $\MCKM$.
Bonferroni correction for three tests at $\alpha=0.05$ requires
$p<0.017$ per test.
For the CP phase: $p \approx 1.09/360 \approx 0.003$ (point prediction,
uniform null). For $M_Z/M_W$ and $\theta_{13}^\nu$: precision $<0.1\%$
corresponds to $p < 0.002$ under uniform null on a $[1,1.5]$ and
$[0\degr,45\degr]$ range respectively.
All three pass the corrected threshold.


\begin{figure}[htbp]
\centering
\includegraphics[width=\textwidth]{fig3_precision}
\caption{Relative errors of HFG predictions for SM flavor observables
(log scale, sorted by precision).
Tier~1 (green): geometrically motivated, $<1\%$ error.
Tier~2 (blue): census matches, $<1\%$ error.
Tier~3 (red): $1$--$4\%$ error.
Dashed and dotted lines mark the $1\%$ and $0.1\%$ thresholds.}
\label{fig:precision}
\end{figure}

%---------------------------------------------------------------
\section{Selberg Zeta Function and Spectral Gap}
\label{sec:selberg}

\subsection{First Laplace eigenvalues}

We compute $\lambda_1$ for both manifolds via the Selberg zeta function
\begin{equation}
  Z_M(s) = \prod_{\gamma\,\mathrm{prim}}\prod_{m,n\geq 0}
  \Bigl(1-e^{i(m-n)\varphi(\gamma)}
  e^{-(s+m+n)\ell(\gamma)}\Bigr),
  \label{eq:selberg}
\end{equation}
whose zeros on $\mathrm{Re}(s)=1$ correspond to Laplace eigenvalues
$\lambda = s(2-s)$.
Using the first 200 primitive geodesics and scanning the critical line:
\begin{align}
  \lambda_1(\MPMNS) &= 2.480 \quad (r_1 = 1.493), \\
  \lambda_1(\MCKM)  &= 2.821 \quad (r_1 = 1.604).
\end{align}
Both satisfy $\lambda_1 \geq 1/4$, consistent with the Selberg
eigenvalue conjecture.

\subsection{Spectral gap conjecture}

\begin{conjecture}[Spectral gap selection]
\label{conj:gap}
A necessary condition for a compact hyperbolic 3-manifold to
realise SM-scale mixing angles (all mixing angles
$\theta_{ij} \in [2\degr, 50\degr]$) is $\lambda_1 \gtrsim 2$
(equivalently $r_1 \gtrsim 1.4$).
\end{conjecture}

The heuristic motivation: manifolds with small $\lambda_1$ have
long, nearly-flat geodesics whose axes cluster at small separations,
producing mixing angles near $0\degr$ or $90\degr$ rather than SM-scale
values.
The large spectral gap $\lambda_1 \sim 2.5$--$2.8$ restricts the
geodesic geometry to a range that naturally produces intermediate angles.

\subsection{Odd-torsion conjecture}

\begin{conjecture}[Odd-torsion selection]
\label{conj:torsion}
A compact hyperbolic 3-manifold with $H_1$ containing 2-torsion cannot
realise SM-scale mixing matrices via the HFG construction.
All manifolds achieving Frobenius fitness $\mathcal{F}<0.020$ in our
scans have $H_1$ with odd-prime torsion of order in $\{3,5,7,13\}$.
\end{conjecture}

%---------------------------------------------------------------
\section{Predictions and Open Problems}
\label{sec:predictions}

\paragraph{Falsifiable predictions.}
\begin{enumerate}[(1)]
\item $\delta_{\mathrm{PMNS}} = 195.91\degr$, testable by
  Hyper-Kamiokande and DUNE ($\sim5\degr$ precision).
\item $M_Z/M_W = 1.1348$ (geometric, from $\MCKM$ geodesics);
  current PDG uncertainty $<0.001\%$, so this is a postdiction
  rather than a prediction.
\item A word of length $\ell^* \approx 12$--$14$ on $\MCKM$ should
  give $|\varphi| \approx 0.20\degr$, matching $\theta_{13}^q$.
\item The prime $13$ should not appear in algebraic covers of $\MPMNS$
  at any degree (testable by degree-12 computation).
\end{enumerate}

\paragraph{Open problems.}
\begin{enumerate}[(1)]
\item Explain $J_\mathrm{CKM} = 0$ at zeroth order and the
  correction mechanism for $J_\mathrm{exp} \approx 3\times10^{-5}$.
\item Prove Conjectures~\ref{conj:gap} and~\ref{conj:torsion}
  rigorously.
\item Derive $\sigma_\mathrm{opt} = (3/2)\log\phi$ geometrically
  from the manifold invariants.
\item Explain the volume adjacency $\mathrm{vol}(\MCKM) \approx
  \mathrm{vol}(m003_\mathrm{cusp})$ to $0.05\%$.
\end{enumerate}

%---------------------------------------------------------------
\section{Conclusions}
\label{sec:conclusions}

We have performed a systematic census of twist angles of the holonomy
of two compact arithmetic hyperbolic 3-manifolds and identified
multiple numerical correspondences with SM flavor parameters.
Three results are especially precise and geometrically motivated:
the leptonic CP phase ($0.55\%$), the $M_Z/M_W$ ratio ($0.024\%$),
and the reactor angle $\theta_{13}^\nu$ ($0.07\%$).
The census enrichment factor ($1.24\times$, $p\approx0.04$) is modest
and should be interpreted cautiously; the Tier-1 results rest on
geometric motivation beyond the census statistics.
The large spectral gap $\lambda_1 \sim 2.5$--$2.8$ suggests a
geometric selection mechanism for SM-scale mixing angles.

%---------------------------------------------------------------
\section*{Data Availability}
All scripts at \url{https://github.com/drmlgentry/hyperbolic-flavor-scan}.
Canonical reproduction:
\texttt{conda run -n sage python reproduce/twist\_reproduce.py}

\section*{Acknowledgments}
The author thanks the developers of SnapPy and SageMath.
During preparation of this manuscript, the author used Claude
(claude-sonnet-4-6, Anthropic) and DeepSeek (DeepSeek-V3) as
research and writing assistants.
All results were independently verified computationally.

\begin{thebibliography}{99}

\bibitem{PDG2024}
S.~Navas et al.\ (Particle Data Group),
Phys.\ Rev.\ D \textbf{110}, 030001 (2024).

\bibitem{SnapPy}
M.~Culler, N.~M.~Dunfield, M.~Goerner, J.~R.~Weeks,
SnapPy, \url{http://snappy.computop.org}.

\bibitem{GMM}
D.~Gabai, R.~Meyerhoff, P.~Milley,
J.\ Amer.\ Math.\ Soc.\ \textbf{22}, 1157 (2009).

\bibitem{MaclachlanReid}
C.~Maclachlan, A.~W.~Reid,
\emph{The Arithmetic of Hyperbolic 3-Manifolds}, Springer (2003).

\bibitem{Selberg56}
A.~Selberg,
J.\ Indian Math.\ Soc.\ \textbf{20}, 47 (1956).

\bibitem{Gentry:Unified}
M.~L.~Gentry,
``Hyperbolic Flavor Geometry: Mixing, CP Violation, and Fermion Masses
from Arithmetic Hyperbolic 3-Manifolds,''
SSRN \href{https://doi.org/10.2139/ssrn.6775158}{6775158} (2026).

\bibitem{Gentry:CKM}
M.~L.~Gentry,
``Quark Mixing from Hyperbolic Geometry,''
SSRN \href{https://doi.org/10.2139/ssrn.6583550}{6583550} (2026);
submitted to J.\ Geom.\ Phys.\ (JGP13076).

\bibitem{Gentry:PMNS}
M.~L.~Gentry,
``Lepton Mixing from Borel Structure of Hyperbolic Holonomy,''
SSRN \href{https://doi.org/10.2139/ssrn.6583549}{6583549} (2026).

\bibitem{Gentry:CP}
M.~L.~Gentry,
``CP Violation from Twist Angles of Hyperbolic Holonomy,''
SSRN \href{https://doi.org/10.2139/ssrn.6583551}{6583551} (2026).

\bibitem{Gentry:FareyTower}
M.~L.~Gentry,
``A Quadratic Cover Prime Formula for a Farey Tower of Arithmetic
Hyperbolic 3-Manifolds,''
SSRN \href{https://doi.org/10.2139/ssrn.6808878}{6808878} (2026);
submitted to Experimental Mathematics (ID: 266619148).

\bibitem{GitHub}
M.~L.~Gentry,
\url{https://github.com/drmlgentry/hyperbolic-flavor-scan}.

\end{thebibliography}

\end{document}
